\documentclass[preprint,12pt,authoryear]{elsarticle}

\usepackage[UTF8,scheme=plain]{ctex}
\usepackage{amsmath,amssymb,mathtools}
\usepackage{booktabs,multirow,array}
\usepackage{graphicx}
\usepackage{xcolor}
\usepackage{enumitem}
\usepackage{lineno}
\usepackage{url}

\journal{Computers \& Industrial Engineering（中文审查版）}
\modulolinenumbers[5]

\definecolor{tbdred}{RGB}{170,0,0}
\newcommand{\TBD}[1]{\textcolor{tbdred}{[待补：#1]}}
\newcommand{\method}{RDMCT-HPS}
\newcommand{\project}{RDMCT-A3C}
\newcommand{\cmax}{C_{\max}}
\newcommand{\pdi}{\mathrm{PDI}}
\newcommand{\gap}{\mathrm{Gap}}

\begin{document}

\begin{frontmatter}

\title{双源混流旋转腔室组合设备的精确MIP调度、两阶段稳定性改进与热启动层次策略搜索}

\author[inst1]{\TBD{第一作者}}
\author[inst1]{\TBD{共同作者}}
\author[inst2]{\TBD{通信作者}\corref{cor1}}
\ead{\TBD{通信作者邮箱}}
\cortext[cor1]{通信作者。}

\affiliation[inst1]{organization={\TBD{院系和学校}},
                    city={\TBD{城市}},
                    postcode={\TBD{邮编}},
                    country={\TBD{国家}}}
\affiliation[inst2]{organization={\TBD{院系和学校}},
                    city={\TBD{城市}},
                    postcode={\TBD{邮编}},
                    country={\TBD{国家}}}

\begin{abstract}
双源混流半导体组合设备在共享旋转腔室内耦合产品晶圆和可复用工艺环境载具，
并受到批处理、清洗、搬运和驻留约束。本文首先将有限时域调度问题建立为精确的
紧凑混合整数规划（MIP）。最大完工时间优先的两阶段路线利用剩余预算保持当前
最佳incumbent的最大完工时间、固定其大部分离散决策，并降低线性等待--节拍
稳定性指标。只有第一阶段被证明最优时才声称主目标最优；第一阶段超时后的改进
不被表述为全局字典序证明。相较于四指机械手四空间工艺模块已有的稳态、单产品、
固定序列线性规划，本文研究不同的有限时域系统，并在若干维度加入更丰富特征。
随后，
本文将结构保持二元骨架（SPBS）原始热启动、层次割选择策略、束搜索解码与SCIP
分支割求解器集成。该策略继承HEM的选择数量/有序子集分解，并增加变量角色
描述量和确定性的角色次模补全。学习回调只排序求解器产生的合法不等式，因而
不改变调度问题的可行域。冻结的配对协议分别检验热启动、层次策略、束搜索和
结构特征在标称及分布外实例上的作用。\TBD{待新检查点和全部冻结实验通过来源
审计与独立解验证后，填入经验证的主实验、组件消融和OOD结果。} 本框架在不把
HEM已有组件误列为本文创新的前提下，将可审计制造物理与学习辅助精确优化结合。
\end{abstract}

\begin{highlights}
\item 精确MIP描述双源混流旋转腔室有限时域调度。
\item 模型增加清洗、可复用载具和有限时域决策。
\item 第二阶段保持incumbent最大完工时间并降低不稳定性。
\item 所有求解方法共享同一SPBS可行解热启动。
\item 在HEM式割选择基础上检验束搜索和变量角色结构。
\end{highlights}

\begin{keyword}
半导体制造 \sep 集群设备调度 \sep 混合整数规划 \sep
割平面选择 \sep 强化学习 \sep 次模优化
\end{keyword}

\end{frontmatter}

\linenumbers

\section{引言}
\label{sec:introduction-zh}

集群设备将多个工艺模块和机械手集成在一个紧密耦合的搬运系统周围。它能够
减少占地面积和污染，但其调度必须同时协调加工、搬运、腔室、装载锁、载具
可用性及配方相关的时间限制。双源混流场景进一步增加了难度：产品晶圆和可复用
工艺环境载具从不同源头进入设备，在共享腔室中汇合，并沿相互作用但不同的
路径流动。批处理和成对处理模式还会形成影响整个调度的离散决策。

精确混合整数规划具有能够证明调度质量、透明表达物理约束等优点，但排序和
资源冲突析取会产生巨大的分支割搜索空间。通用求解器在一个节点可能生成大量
候选割。全部加入虽然可能强化松弛，却也会增大线性规划规模、数值负担和回调
开销。因此，割选择是一个序贯决策问题，局部割效能无法充分描述割之间的相互
作用。

已有研究表明，机器学习可以在不替代数学模型的前提下辅助组合优化
\citep{bengio2021machine}。强化学习已被用于割选择\citep{tang2020cut}，前瞻
模仿学习则提供另一条路线\citep{paulus2022looking}。尤其是HEM已经用指针网络
同时学习所选割比例和有序子集，
并采用与A3C密切相关的并行策略梯度训练。因此，层次数量预测、有序子集、指针
网络和并行采样在本文中均属于已有基线，而不是新贡献。本文考察的是：在HEM
嵌入精确调度求解器后，模型导出的原始热启动、有限宽束搜索及模型角色信息能否
改善其随时性能。

本文研究两个相互关联的问题。第一，有限时域精确模型能否在稳态固定序列范围
之外，统一表达双物料源、混合$4\times1$/$2\times2$流、旋转腔室、清洗和
可复用载具？第二，SPBS初始化和束搜索能否在不改变模型有效性的前提下改善
HEM式学习割选择器的随时性能？为此，本文建立精确紧凑模型和集成求解框架
\method{}。模型把已由制造等式固定的变量连续表示，但不松弛真正的分配、排序、
模式或清洗决策；求解栈组合共享SPBS可行解、SCIP、23维角色感知层次策略和
有限宽束搜索。

本文主要贡献如下。

\begin{enumerate}[leftmargin=*]
  \item 将含共享机械手、旋转腔室、两种配方模式、清洗、载具驻留和可复用令牌
  的有限时域双源混流问题建立为精确紧凑MIP。状态感知的第二阶段保持当前最佳
  incumbent最大完工时间，并在大部分离散决策固定时改进线性等待--节拍指标。
  \item 构造SPBS热启动：固定具有物理意义的分配、激活和清洗二元骨架，再由
  SCIP补全可行解。配对比较中所有方法共享同一可行解及其构造时间。
  \item 相对适配HEM基线保留其层次数量和有序指针策略，进一步检验有限宽束搜索、
  10个变量角色描述量及透明的角色次模补全；创新主张限定为上述集成，并
  必须由组件级消融支持。
  \item 建立冻结的实例配对协议，覆盖模型验证、热启动与搜索消融、标称比较、
  析因条件、SPBS初始化、单因素敏感性以及规模/配方分布偏移。
\end{enumerate}

当前文档是一份遵循证据纪律的初稿：在所有必需配对实验完成之前，正式数值
单元格保持为空，以避免选择性报告不完整种子或仅报告成功求解实例。

\section{相关研究}
\label{sec:related-zh}

\subsection{半导体集群设备调度}

集群设备调度把机器调度与机械手搬运、驻留时间约束相结合。已有研究针对包含
集群设备的光刻区域\citep{madathil2018photolithography}以及带清洗操作的
时间受限设备\citep{yuan2023purge}提出了精确或解析方法。与本文最接近的是
\citet{gu2024multifinger}：该研究面向四指机械手与四空间工艺模块，在两个预先
给定的循环机械手序列---后向序列（BS）和混合任务序列（HTS）---下分析稳态
运行，并分别求解一个满足驻留约束、以周期时间最小为目标的线性规划。该文采用
单一产品流，并把腔室清洗和多种晶圆同时加工留作后续研究。本文采用不同的搬运
架构，并在若干维度加入更丰富特征，包括有限时域、两个物料源、产品/载具混流、
灵活分配与排序、共享旋转腔室、两种配方模式、显式清洗及可复用载具令牌。目前
尚未证明本文模型能够精确退化为其四指洁净/污染手指系统。因此，本文比较建模
覆盖范围和刻意构造特例，而不比较两个非等价模型的原始运行时间。

\subsection{精确MIP中的机器学习}

SCIP提供模块化分支--割--定价框架，使学习组件能够在受控决策点介入
\citep{achterberg2009scip,bestuzheva2023scip}。当学习组件只改变搜索优先级
而不改变数学有效性时，学习辅助精确优化更容易审查和部署。基于这一原则，
本文回调只接收SCIP生成的合法割，选择其中的优先子集，并让未选割继续供求解器
使用；学习模型既不生成不等式，也不剪枝。

\subsection{学习式割选择}

\citet{tang2020cut}将割选择构造为强化学习问题，\citet{paulus2022looking}
利用前瞻信息学习能带来有效界改进的割。\citet{wang2023hem}使用13个通用割
特征和指针网络\citep{vinyals2015pointer}，层次建模所选割数量与顺序；其并行
层次策略梯度被描述为与A3C密切相关。本文不把层次数量预测、有序子集选择、
指针网络或并行采样作为创新，而保留该基线并检验三项增加：策略外部的SPBS
原始热启动、测试时有限宽束搜索，以及带确定性次模补全的模型角色特征。当前
正式配置采用移动平均REINFORCE基线，因此论文将训练算法称为层次策略梯度；
仓库名\project{}本身不能支持严格异步优势演员--评论家（A3C）的主张。

\section{问题描述与精确紧凑模型}
\label{sec:model-zh}

\subsection{制造系统拓扑}

产品晶圆的路径为
\[
 \mathrm{LP}\rightarrow\mathrm{ATR}\rightarrow\mathrm{AL}
 \rightarrow\mathrm{ATR}\rightarrow\mathrm{LL_U}
 \rightarrow\mathrm{VTR}\rightarrow\mathrm{CH}
 \rightarrow\mathrm{VTR}\rightarrow\mathrm{LL_L}
 \rightarrow\mathrm{ATR}\rightarrow\mathrm{LP},
\]
其中，LP为装载端口，ATR和VTR分别为大气机械手和真空机械手，AL为校准器，
$\mathrm{LL_U}$与$\mathrm{LL_L}$分别为上、下装载锁。可复用工艺环境载具
（PEC）经VTR在PEC存储区和共享腔室之间流动。CH2和CH3为共享搬运和装载资源
的旋转腔室。

系统包含两种配方模式。在$4\times1$模式中，四片产品晶圆使用两个可达槽位组
之一进行批处理；在$2\times2$模式中，一个PEC和两片产品晶圆构成物理有序的
PEC--产品--产品--PEC链。清洗可在满足清洗间隔约束时提前触发；所选PEC必须
在相应链持续期间保持驻留，并通过可复用令牌建模。资源锁禁止机械手动作和
不相容装载操作并发执行。

表~\ref{tab:model-scope-zh}明确本文相对最接近四空间工艺模块研究的主张边界。
表中的内容表示模型显式优化的决策，而不是设备可能具备但未纳入优化的物理能力。

\begin{table}[htbp]
\centering
\caption{相对\citet{gu2024multifinger}的建模范围。}
\label{tab:model-scope-zh}
\begin{tabular}{p{0.25\linewidth}p{0.31\linewidth}p{0.34\linewidth}}
\toprule
维度 & Gu等（2024） & 本文精确模型 \\
\midrule
时域 & 稳态周期 & 有限时域事件调度 \\
搬运架构 & 四指洁净/污染手指机械手 & 分离的大气/真空搬运机械手 \\
物料流 & 单一产品流 & 产品晶圆与可复用PEC \\
配方 & 每个算例一种流模式 & 混合$4\times1$与$2\times2$模式 \\
机械手决策 & 预先固定BS或HTS & 优化分配和资源次序 \\
旋转 & 四空间工艺模块旋转延迟 & 共享旋转腔室/槽位组事件 \\
清洗 & 留作未来研究 & 清洗周期和提前清洗决策 \\
数学规划 & 两个序列专用LP & 一个精确混合整数模型 \\
目标 & 稳态周期时间 & 最大完工时间；条件式incumbent保持稳定性改进 \\
\bottomrule
\end{tabular}
\end{table}

\subsection{事件和资源变量}

设$\mathcal{O}$为操作集合，$\mathcal{R}$为物理资源集合，$\mathcal{A}$为
工艺路径产生的优先弧集合。操作$i\in\mathcal{O}$的开始时间、加工时间和完成
时间分别为$s_i$、$p_i$和$c_i=s_i+p_i$。二元变量$y_{ijr}$确定在单位容量
资源$r$上相互冲突的两个操作的先后顺序；$x$表示分配及模式选择变量，$u$
表示PEC令牌分配变量。模型包含优先约束
\begin{equation}
  s_j \ge c_i,\qquad (i,j)\in\mathcal{A},
  \label{eq:precedence-zh}
\end{equation}
以及资源析取约束
\begin{align}
  s_j &\ge c_i-M_{ijr}(1-y_{ijr}), \\
  s_i &\ge c_j-M_{jir}y_{ijr},
  \qquad \{i,j\}\subseteq\mathcal{O}_r,
  \label{eq:disjunction-zh}
\end{align}
其中$M_{ijr}$为事件相关的有效常数。容量、腔室旋转、槽位组、清洗、驻留和
令牌守恒约束将$x$、$u$、$y$与$s$连接起来。完整的索引级模型和参数构造将在
\TBD{在线补充材料及代码仓库链接}中给出。

\subsection{等价紧凑化}

部分名义上的二元成员、槽位和别名变量已经被所选制造模式产生的等式固定为
0或1。紧凑模型在保留这些固定等式的前提下，将相应变量声明为连续变量。这一
操作既不改变变量取值，也不改变可行域的投影，只是去除多余的整数声明。所有
真正的分配、排序、模式和清洗决策仍保持为整数。因此，本文的“紧凑”表示等价
重构，而不是对制造物理约束的松弛。

\subsection{最大完工时间优先的两阶段改进}

主目标为第一片产品晶圆离开LP到最后一片产品晶圆返回LP的时间区间。第一阶段在
分配预算内尝试求解精确MIP
\begin{equation}
  \min_{z\in\mathcal{F}} \cmax(z).
  \label{eq:makespan-zh}
\end{equation}
记其最佳incumbent为$\widehat z$，$\widehat C=\cmax(\widehat z)$，求解状态为
$\sigma_1$。若无incumbent，则跳过第二阶段；否则固定大部分离散分配/排序决策
$d\in\mathcal{D}_{\mathrm{fix}}$，保留部分双夹持决策自由，并加入
\begin{equation}
  z_d=\widehat z_d\quad(d\in\mathcal{D}_{\mathrm{fix}}),
  \qquad \cmax(z)\le\widehat C+\epsilon.
  \label{eq:incumbent-fix-zh}
\end{equation}
冻结的线性稳定性目标为
\begin{equation}
 \min S_{\mathrm{lin}}(z)=
 \alpha_w\sum_{i\in\mathcal{W}}w_i
 +\alpha_m w^{\max}
 +\alpha_c\sum_{j\in\mathcal{C}}\delta_j
 +\alpha_e\sum_{k\in\mathcal{E}}e_k,
 \label{eq:stability-linear-zh}
\end{equation}
其中$w_i$为等待/空闲量，$w^{\max}$为最大单项等待，$\delta_j$为相邻启动节拍
差的绝对偏差，$e_k$为软等待上限超额量；所有权重在评估前固定。确认性配置为
该阶段预留60秒和至多5000个节点，并分别记录$\sigma_1$及第二阶段状态。只有
$\sigma_1=\texttt{optimal}$时，$\widehat C$才是已证明的主目标最优值；主阶段
超时时它只是当前最佳incumbent。由于大部分离散决策已固定，第二阶段属于保持
incumbent的调度抛光，而不是全局字典序稳定性最优证明。

\section{热启动层次割选择与精确求解}
\label{sec:method-zh}

\subsection{集成求解栈与SPBS热启动}

主求解路线为
\[
 \text{SPBS可行解}\rightarrow
 \text{HEM式策略/束搜索}\rightarrow
 \text{SCIP分支割}\rightarrow
 \text{带分阶段状态的两阶段调度}.
\]
SPBS固定具有物理意义的二元骨架，包括批次与腔室分配、激活指示、清洗周期、
槽位选择和部分资源次序，再由一个以可行性为目标的短时SCIP过程补全连续事件
时间及有意保留的未固定决策。只有SCIP判定可行时，该解才作为主MIP的原始热
启动。热启动构造属于预处理，其墙钟时间、成功标志和解哈希均需记录。每个配对
对比中的所有方法接收同一个解文件；若SPBS在预算内失败，则确认性配对标记为
不完整并重新生成热启动，不能静默降级为无热启动观测。

SPBS和学习割选择器针对不同瓶颈：前者主要缩短首次可行解时间，后者排序求解器
产生的割以改善原始--对偶轨迹。因此，二者的独立效应与交互效应必须分别估计。

\subsection{决策接口与学习目标}

在一次分离轮次中，SCIP提供合法候选池
$\mathcal{P}=\{1,\ldots,n\}$，候选$i$表示有效不等式
$a_i^\top z\le b_i$。低成本效能预筛选最多保留128个候选，策略最多选择16个
割，未选候选按原合法顺序归还。

状态包含候选特征和求解器上下文，动作由选择数量和有序子集组成。训练奖励为
冻结SCIP版本报告的原始--对偶积分负值：
\begin{equation}
 R=-\pdi_{\mathrm{SCIP}},
 \label{eq:pdi-zh}
\end{equation}
PDI越低，随时性能越好。最终补充材料将明确冻结SCIP版本的界归一化、初始gap及
无incumbent约定；本文不假定它与HEM实现数值口径完全一致。
对于基线$b$，策略梯度估计为
\begin{equation}
 \widehat{\nabla_\theta J}
 =\frac{1}{B}\sum_{\ell=1}^{B}
   (R_\ell-b_\ell)\nabla_\theta
   \log\pi_\theta(a_\ell\mid s_\ell),
 \label{eq:reinforce-zh}
\end{equation}
该估计遵循\citet{williams1992reinforce}。独立训练样本由并行工作进程求值，
并在更新前聚合。冻结配置使用移动平均奖励基线，而不是学习得到的评论家。因此，
本文将该实现准确表述为并行层次策略梯度，而不是严格异步优势演员--评论家。
仓库名\project{}仅为保持可追溯性；只有在运行独立演员--评论家实现及预先指定的
对比之后，才能在算法层面直接声称A3C。

\subsection{通用特征与模型角色特征}

13维通用向量包含目标平行度、效能、支撑度、整数支撑度、归一化违反程度，
以及割系数和目标系数的均值、最大值、最小值与标准差。

结构扩展仅根据SCIP变量类型及目标系数，把变量划分为五种通用类型/角色：二元、
一般整数、目标系数非零的连续变量、其他连续辅助变量及无法分类变量。它不根据
变量名推断分配、清洗或腔室语义。对割$i$和角色$k$，其系数质量占比为
\begin{equation}
 m_{ik}=
 \frac{\sum_{j\in\mathcal{V}_k}|a_{ij}|}
      {\sum_j|a_{ij}|+\epsilon}.
 \label{eq:mass-zh}
\end{equation}
五个占比再加上有界变量比例、正系数比例、离散--连续耦合度、主导角色占比和
归一化角色熵，形成23维向量。新增10个比例量在数值容差内对约束行正比例缩放
不变，并依赖变量类型而非具体实例的变量名；完整23维向量不声称尺度不变，因为
冻结的未归一化13维基线包含原始系数统计量。

\subsection{层次策略与束搜索}

数量策略首先预测选择比例，随后指针式策略对有序序列评分
\citep{vinyals2015pointer}。网络使用128维隐藏表示和一次glimpse。训练时
采用随机采样，评估时采用贪心解码或宽度为3的束搜索。各学习基线共享序列结束
处理、累计对数概率、回溯指针和候选池，确保束搜索比较不会改变候选范围或
计时口径。HEM本身已经包含数量策略和有序指针策略；本文作为搜索扩展检验的
只是用有限宽束搜索替代贪心推理。束搜索推理和回调时间均计入共同墙钟上限。

\subsection{角色次模补全}

设$E$为由策略排序及启发式尾部截断得到的扩展池，$A\subseteq E$为固定策略锚点，
$S\supseteq A$为补全集合。令$q_i$为归一化质量，
$r_i=1-\operatorname{rank}(i)/\max(1,|\mathcal{R}|-1)$为排序先验，$m_{ik}$为
角色质量，$s_{ij}\in[0,1]$为实现特征的截断余弦相似度，并令归一化需求权重
$d_k\propto\sum_{i\in E}m_{ik}$、$\sum_kd_k=1$。补全过程求解
\begin{align}
F_E(S)=&\;w_q\sum_{i\in S}q_i+w_p\sum_{i\in S}r_i
 +w_c\sum_k d_k\sqrt{\sum_{i\in S}m_{ik}} \nonumber\\
&+\frac{w_r}{|E|}
 \sum_{j\in E}\max_{i\in S}s_{ij},
\label{eq:submodular-zh}
\end{align}
且满足$A\subseteq S\subseteq E$、$|S|\le K$。归一化默认权重为
$(w_q,w_p,w_c,w_r)=(0.35,0.20,0.25,0.20)$，补全池系数为2.0，锚点比例为
0.25。

在固定$E$上的非负模函数项、需求加权凹角色覆盖项及设施选址代表性项下，$F_E$
为单调次模函数。给定锚点和扩展池后，贪心剩余补全满足
\begin{equation}
 F_E(A\cup S_g)-F_E(A)\ge(1-1/e)
 \max_{|B|\le K-|A|}\{F_E(A\cup B)-F_E(A)\},
\end{equation}
其中$B\subseteq E\setminus A$\citep{nemhauser1978analysis}。该保证不覆盖锚点
选择、扩展池截断、分支割运行时间或最终MIP质量。

\subsection{合法性与开销}

\method{}不生成约束、不修改系数，也不宣布节点不可行。每个被选不等式均已
通过SCIP的合法性检查，未选割仍按回调契约供求解器使用。因此，学习组件能够
改变搜索轨迹和运行时间，但不改变调度模型的可行集，也不影响最终证明的有效
性。所有回调推理、束搜索和后处理时间均计入统一的墙钟时间上限。

\subsection{可选RL-SAT路线的证据边界}

代码库还包含可选的A3C--束搜索--CP-SAT路线。该路线搜索腔室分配，并针对每个
候选求解一个CP-SAT时间抽象；其中的“optimal”只证明该候选专用抽象达到最优，
不能证明Petri网MIP或全部腔室分配的全局最优。因此，RL-SAT不纳入主比较。只有
在完成特征同等性、独立调度验证、统一端到端计时及证书范围检查后，才可作为
独立附录结果报告。

\section{实验设计}
\label{sec:experiments-zh}

\subsection{冻结测试集}

表~\ref{tab:suites-zh}汇总冻结测试集。标称实例用于主要比较；析因测试集改变
晶圆数、配方构成和加工时间尺度；敏感性测试一次改变一个物理因素；OOD测试
同时增大规模并改变配方平衡。测试实例不用于超参数调优。

\begin{table}[htbp]
\centering
\caption{冻结评估测试集。}
\label{tab:suites-zh}
\begin{tabular}{p{0.18\linewidth}p{0.55\linewidth}r}
\toprule
测试集 & 设计 & 条件/实例数 \\
\midrule
标称主实验 & 冻结的标称制造实例 & 20 \\
析因DOE & 晶圆数$\{8,16,24,32\}$、五种配方比例、
加工尺度$\{0.9,1.0,1.1\}$ & 60 \\
SPBS & 符合条件的core实例，无原始解与同一自动生成SPBS解 & 48对 \\
敏感性 & 搬运$\{0.8,1.2\}$、加工$\{0.8,1.2\}$、
PEC容量10、清洗间隔$\{3,10\}$ & 8 \\
OOD & 晶圆数$\{40,48,64\}$，配方比例
$\{3{:}1,1{:}1,1{:}3\}$ & 9 \\
\bottomrule
\end{tabular}
\end{table}

\subsection{比较方法}

\begin{table}[htbp]
\centering
\caption{冻结主实验中的方法。所有方法使用相同模型、实例、热启动文件、资源
上限、求解器种子和全局SCIP设置。学习选择器另共享效能预筛与选择上限；SCIP
和ACS保留各自原生候选机制。}
\label{tab:methods-zh}
\begin{tabular}{p{0.22\linewidth}p{0.69\linewidth}}
\toprule
方法 & 说明 \\
\midrule
SCIP & SCIP默认割选择。 \\
ACS & 仅在冻结验证集上调参的确定性自适应割选择基线。 \\
HEM-13D（适配） & 13维通用特征、层次策略、贪心解码，无变量角色补全。 \\
HEM-13D+beam & 使用宽度3束搜索的适配HEM。 \\
23D feature only & 通用特征加角色特征，不使用角色次模重排。 \\
Structure+greedy & 23维策略、角色次模补全和贪心解码。 \\
\method{} & 23维策略、宽度3束搜索、策略锚点和角色次模补全。 \\
\bottomrule
\end{tabular}
\end{table}

适配HEM-13D保留\citet{wang2023hem}的选择数量和有序指针思想，但采用本文的
效能top-128预筛、16割上限、训练预算、SCIP版本和调度协议，不能称为逐比特复现。
HEM-13D+beam只改变测试时解码。23D和Structure两行依次隔离变量角色比例和
确定性补全。另设热启动实验，在SCIP、HEM-13D和\method{}下分别比较无热启动与
同一SPBS文件；热启动--策略交互采用差分中的差分，而不是单个成对差。完成该
矩阵前，不能把改进归因于集成的SPBS+策略+束搜索栈。\TBD{确认性benchmark前
实现并审计学习策略none/auto Stage；当前SPBS Stage仅覆盖SCIP。}

\subsection{训练与实现}

冻结策略梯度配置采用5个独立训练种子、60个epoch、低层策略$10^{-4}$的Adam
学习率、批量大小8、最大梯度范数2、每个epoch 8个训练样本和2个并行样本任务。
高层选择百分比策略的学习率为$5\times10^{-4}$，奖励基线为指数移动平均。
候选数和选择数上限分别为128和16。由于当前训练文件把在线评估频率设为0，
预先指定使用最终检查点（epoch 60）；论文不声称按验证集选择检查点。
\TBD{若声称严格A3C，使用经审计的演员--评论家实现重新训练，并增加
价值损失、更新时序和对应消融。}

每个训练求解限制为30秒和1024 MB，只在根节点进行一轮分离；学习率每10个epoch
乘以0.96。评估前冻结上述设置、训练实例hash和15个选定\texttt{itr\_60.pkl}
的精确hash。

实验通过PySCIPOpt调用SCIP。主实验、DOE、SPBS和敏感性实验的时间上限为
600秒，OOD实验的时间上限为1200秒，并统一采用4096 MB的SCIP内存上限。
所有实验完成后，最终稿将报告SCIP、PySCIPOpt、Python、PyTorch、CUDA、GPU、
CPU、操作系统、线程数和分阶段内存上限。\TBD{补充最终复现环境表以及代码
仓库/归档DOI。}

\subsection{指标与统计分析}

主要指标为PDI。次要指标包括最终原始界、对偶界、相对间隙、incumbent可用率、
分阶段状态和时间、节点数、超时率、内存失败率和最优求解率。\TBD{在确认性
冻结前，增加并审计首次/最佳可行解时间、根节点gap、LP迭代、接受割、所选角色
覆盖、冗余度和回调开销遥测，或删除相应机理分析。}

运行按制造实例、求解器种子及学习方法的训练种子匹配。推断时先按预声明规则在
每个制造实例内聚合求解器/训练重复，以制造实例为抽样单元。对每个预先指定的
对比，报告配对中位数差、bootstrap置信区间、适用时的
Wilcoxon符号秩检验以及Holm校正$p$值。超时和失败必须保留并单独汇总，不得
为了改善表面性能而删除。

求解器状态\texttt{timelimit}本身并不是无效观测。它表示在共同墙钟预算内尚未
完成所需证明；若存在可行 incumbent，该调度仍可能具有运行价值。当测试集中有
未完成实例时，各方法相同的封顶求解时间不能用于排序。按照\citet{wang2023hem}
的处理，主比较使用PDI，并同时报告可行解可用率、最终间隙、上下界和超时率。
无可行解的行保留。\TBD{确认性冻结前，明确并测试最终投稿分析所采用的SCIP版本
初始gap/无incumbent约定，不得引入未记录惩罚。}

\section{实验结果}
\label{sec:results-zh}

\subsection{归档模型健全性检查（非验证性）}

表~\ref{tab:legacy-sanity-zh}记录本地\texttt{\_results}归档中保留的求解器结果。
它们说明若干纯配方及产品--载具混流模型能够产生可行解，并在这四个单例上由
求解器报告零最终间隙。但这些结果早于冻结实验，没有当前要求的独立解验证清单，
且实例参数不一致。因此，它们只作为模型示例，不参与任何统计或算法优势主张。
纳入规则在读取方法时间前固定：每种不同归档纯配方或产品--载具配置选一条可解析、
零gap且$C_{\max}$/节点/时间字段完整的记录。\TBD{在补充材料发布含排除、超时、
错误和重复行的完整归档状态清单。}

\begin{table}[htbp]
\centering
\caption{归档单实例模型检查；数值仅作描述，不属于验证性实验。}
\label{tab:legacy-sanity-zh}
\begin{tabular}{lrrrrr}
\toprule
配置 & 产品晶圆数 & 状态 & $C_{\max}$ & 节点数 & 时间（秒） \\
\midrule
$2\times2$令牌冒烟测试 & 4 & optimal & 285.007 & 972 & 1 \\
纯$4\times1$ & 17 & optimal & 1002.000 & 1 & 3 \\
纯$2\times2$ & 17 & optimal & 1244.000 & 211276 & 549 \\
$4\times1$归档实例 & 33 & optimal & 984.729 & 69 & 9 \\
\bottomrule
\end{tabular}
\end{table}

另一个归档实例包含四个旧算法标签。四种方法均报告同一个已证明的
$C_{\max}=2234$，墙钟时间分别为91秒（SCIP）、89秒（仅A3C）、125秒（启发式
束搜索）和127秒（A3C+束搜索）。这一单例既不能证明加速，也不能证明束搜索有益；
其方向反而说明必须完成配对束搜索消融。只挑选89秒这一行作为证据是无效的。

\subsection{标称主实验}

\begin{table}[htbp]
\centering
\caption{标称主实验结果。结果应由20个实例、5个求解器种子以及学习方法的
5个训练种子计算。PDI、间隙和时间越低越好，最优率越高越好。}
\label{tab:main-zh}
\begin{tabular}{lrrrrr}
\toprule
方法 & PDI & 最终间隙（\%） & 首次可行解（秒） & 节点数 & 最优率（\%）\\
\midrule
SCIP & -- & -- & -- & -- & -- \\
ACS & -- & -- & -- & -- & -- \\
HEM-13D（适配） & -- & -- & -- & -- & -- \\
HEM-13D+beam & -- & -- & -- & -- & -- \\
23D feature only & -- & -- & -- & -- & -- \\
Structure+greedy & -- & -- & -- & -- & -- \\
\method{} & -- & -- & -- & -- & -- \\
\bottomrule
\end{tabular}
\end{table}

\TBD{报告完整配对的标称比较、效应量、置信区间、校正检验和失败数量。不得
使用当前不完整分片填表。}

\subsection{消融与机理分析}

\begin{table}[htbp]
\centering
\caption{预先指定的机理对比。}
\label{tab:ablation-zh}
\begin{tabular}{p{0.31\linewidth}p{0.38\linewidth}rr}
\toprule
对比 & 隔离的机制 & $\Delta$PDI & Holm $p$ \\
\midrule
SCIP+SPBS vs. SCIP（无热启动） & 热启动可行解 & -- & -- \\
HEM-13D+SPBS vs. HEM-13D（无热启动） & HEM条件下热启动效应 & -- & -- \\
$\Delta_{\mathrm{HEM,warm}}-\Delta_{\mathrm{SCIP,warm}}$ & 热启动--策略交互（DiD） & -- & -- \\
HEM-13D+beam vs.\ HEM-13D & 搜索宽度 & -- & -- \\
23D feature only vs.\ HEM-13D & 变量角色特征（双方均greedy） & -- & -- \\
Structure+greedy vs.\ 23D feature only & 次模补全 & -- & -- \\
\method{} vs.\ Structure+greedy & 固定结构规则下的束搜索 & -- & -- \\
\bottomrule
\end{tabular}
\end{table}

\TBD{增加割池大小、所选割角色覆盖、冗余度、回调开销和根节点界分析，将求解
表现与所提机理联系起来。}

\subsection{析因、SPBS和敏感性结果}

\begin{table}[htbp]
\centering
\caption{设计制造条件下的稳健性汇总。}
\label{tab:robustness-zh}
\begin{tabular}{lrrrr}
\toprule
研究 & 条件数 & 主效应/对比 & $\Delta$PDI & 校正$p$\\
\midrule
析因DOE & 60 & \TBD{指定预设模型项} & -- & -- \\
SPBS & 48对 & 自动SPBS原始解 vs.\ 无热启动 & -- & -- \\
敏感性 & 8 & 最大单因素退化 & -- & -- \\
\bottomrule
\end{tabular}
\end{table}

\TBD{只有在交互项已预设且可估计时才报告析因交互。SPBS改变原始初始化而不
改变制造可行域；其主效应和学习策略交互应与割选择效应分开报告。}

\subsection{分布外泛化}

\begin{table}[htbp]
\centering
\caption{按晶圆数划分的OOD结果。所有比较方法共享4096 MB内存上限和1200秒
墙钟时间上限。Inc.表示具有经验证incumbent的运行比例；超时、内存限制和
runner错误率在配套状态表中分别报告。}
\label{tab:ood-zh}
\begin{tabular}{lrrrrrr}
\toprule
\multirow{2}{*}{方法} &
\multicolumn{2}{c}{40片} &
\multicolumn{2}{c}{48片} &
\multicolumn{2}{c}{64片} \\
\cmidrule(lr){2-3}\cmidrule(lr){4-5}\cmidrule(lr){6-7}
& PDI & Inc.（\%） & PDI & Inc.（\%） & PDI & Inc.（\%）\\
\midrule
SCIP & -- & -- & -- & -- & -- & -- \\
ACS & -- & -- & -- & -- & -- & -- \\
HEM-13D（适配） & -- & -- & -- & -- & -- & -- \\
HEM-13D+beam & -- & -- & -- & -- & -- & -- \\
23D feature only & -- & -- & -- & -- & -- & -- \\
Structure+greedy & -- & -- & -- & -- & -- & -- \\
\method{} & -- & -- & -- & -- & -- & -- \\
\bottomrule
\end{tabular}
\end{table}

一项工程活动于2026年8月2日在4096 MB/1200秒设置下完成了计划的729条OOD记录，
且没有运行器错误；但全部记录均以\texttt{timelimit}结束。SCIP和ACS各自的
27条记录全部产生可行解，而每个学习方法族仅在135条中的27条产生可行解：五个
旧训练检查点中只有第一个做到了这一点。此外，五个检查点的元数据均记录
\texttt{experiment.seed=1}，声称的独立训练种子来源无法审计。因此，这729条
记录不进入表~\ref{tab:main-zh}和表~\ref{tab:ood-zh}；完成全部运行并不能修复
检查点来源缺陷。\TBD{用通过验证的新冻结OOD活动统计替换本工程审计段落。}

\section{工业管理启示}
\label{sec:implications-zh}

本框架把三类责任分开：制造工程师定义可审计的物理约束，SCIP负责有效不等式
和最优性证明，学习组件负责分配有限的割预算。当优化器需要得到改进、但可行性
不能交由黑箱预测时，这种职责划分具有实际意义。PDI还刻画了尽早获得高质量
调度的运行价值，即使在派工时限内尚未完成最优性证明。

\TBD{完整实验结束后，将观测到的PDI和可行解时间改善转化为测试晶圆规模下的
生产含义；不得把吞吐量或成本结论外推到未建模系统。}

\section{局限性与有效性威胁}
\label{sec:limitations-zh}

第一，本文模型与\citet{gu2024multifinger}的LP描述非等价设备决策；可以比较
覆盖范围并进行刻意简化的特例交叉检查，但不能直接比较原始运行时间。第二，
精确性主张属于MIP可行域及求解器报告的主阶段证明；第二阶段固定大部分incumbent
离散决策，不是全局字典序证书。第三，角色表示依赖模型变量到模型角色的正确
映射，且尚未在无关MIP上验证。第四，次模保证只适用于剩余集合补全，不是求解
时间保证。第五，当前主训练代码是并行
REINFORCE；若不重新运行演员--评论家路线，严格A3C主张不在现有证据范围内。
第六，墙钟比较对硬件和实现敏感，热启动构造、推理和束搜索均须计时。OOD内存
压力是实验结果，而不是可删失的不便。最后，RL-SAT目前证明的是CP-SAT抽象而非
完整Petri MIP；实际部署还需考虑不确定性、重调度、设备故障和晶圆厂级集成。

\section{结论}
\label{sec:conclusion-zh}

本文面向双源混流旋转腔室调度，给出精确紧凑MIP和集成的热启动层次策略搜索
框架。模型统一表示有限时域产品/PEC流、混合配方模式、清洗、腔室旋转及共享
机械手资源。状态感知的第二阶段保持主阶段incumbent最大完工时间，并改进线性
等待--节拍指标；它不被报告为全局字典序证明。求解路线组合SPBS可行解、已有
HEM数量/有序子集策略、有限宽束搜索、变量角色特征和SCIP。回调只作用于求解器
生成的合法割，因此保持模型有效性。本文已建立冻结配对协议，以分别检验模型、
热启动、策略和束搜索机制。
\TBD{只有在新检查点、独立解验证和全部配对活动通过审计后，填入最终证据结论。}

\section*{CRediT作者贡献声明}

\TBD{填写每位作者的CRediT角色。}

\section*{经费}

\TBD{填写经费信息，或声明未获得专项经费。}

\section*{利益冲突声明}

\TBD{填写经所有作者确认的声明。}

\section*{数据与代码可用性}

仓库保留项目名\project{}；由于冻结训练路线是策略梯度而非严格演员--评论家A3C，
论文使用中性算法名\method{}。冻结配置、实例生成器、模型实现、训练代码和分析脚本将在
\TBD{代码仓库URL和不可变DOI}归档。原始实验输出及软硬件版本的机器可读清单
将随归档一同提供。

\section*{致谢}

\TBD{填写致谢或删除本节。}

\bibliographystyle{elsarticle-harv}
\bibliography{rdmct_cie_references}

\end{document}
